\chapter{Programmazione di rete: i socket in C}

Finora abbiamo visto applicazioni e protocolli esistenti. Ora vediamo come \textbf{creare} un'applicazione
di rete. Quasi tutte sono fatte di un programma lato client e di un programma lato server: quando li
eseguiamo nascono un processo client e un processo server, che comunicano \textbf{leggendo e scrivendo sui
socket}.

% ══════════════════════════════════════════════════════════════
\section{Le scelte di progetto}
% ══════════════════════════════════════════════════════════════

\subsection{Protocolli aperti e proprietari}

\begin{itemize}
    \item \textbf{Protocolli aperti}: le regole sono pubbliche (per esempio in un RFC). Un client scritto da noi
          può parlare con un server scritto da altri, purché entrambi rispettino lo standard (un nostro client HTTP
          parla con Apache).
    \item \textbf{Protocolli proprietari}: regole ``fatte in casa'', non pubblicate. Client e server li scrive lo
          stesso sviluppatore, che ha pieno controllo ma deve fare tutto da sé.
\end{itemize}

\subsection{La scelta del protocollo di trasporto}

\begin{itemize}
    \item \textbf{TCP}: orientato alla connessione, offre un \textbf{canale affidabile a flusso di byte} tra i due end system.
    \item \textbf{UDP}: senza connessione, invia \textbf{pacchetti indipendenti} da un end system all'altro, senza
          alcuna garanzia di consegna.
\end{itemize}

% ══════════════════════════════════════════════════════════════
\section{Dove stanno i socket nella pila}
% ══════════════════════════════════════════════════════════════

\begin{figure}[H]
\centering
\begin{tikzpicture}[every node/.style={font=\small\sffamily}]
  \node[lay app, minimum width=3.6cm, minimum height=0.9cm] (a) at (0,1.8) {Applicazione};
  \node[lay tra, minimum width=3.6cm, minimum height=0.9cm] (t) at (0,0.9) {Trasporto (TCP/UDP)};
  \node[lay net, minimum width=3.6cm, minimum height=0.9cm] (n) at (0,0) {Rete (IP)};
  \node[lay link, minimum width=3.6cm, minimum height=0.9cm] (l) at (0,-0.9) {Accesso (link + fisico)};
  \node[draw, fill=lapp!40, rounded corners, minimum width=3.4cm, minimum height=0.9cm] (my) at (6,1.8) {la vostra applicazione};
  \node[draw, fill=white, very thick, draw=netred, rounded corners, minimum width=3.4cm, minimum height=0.5cm] (api) at (6,1.1) {\textbf{API dei socket}};
  \node[draw, fill=lphy!50, rounded corners, minimum width=3.4cm, minimum height=1.8cm, align=center] (os) at (6,-0.2) {sistema operativo\\+ hardware};
  \draw[dashed, netred, thick] (-2.4,1.35) -- (8.2,1.35);
  \node[text=netred, anchor=west, font=\scriptsize\sffamily] at (8.3,1.35) {socket};
  \draw[decorate, decoration={brace, amplitude=5pt, mirror}] (-2.1,1.3) -- (-2.1,-1.35) node[midway, left=6pt, align=center, font=\scriptsize\sffamily] {implementati\\dal sistema\\operativo};
\end{tikzpicture}
\caption{Dove sono implementati gli strati: l'applicazione è nostra, trasporto e rete stanno nel sistema operativo e si usano attraverso l'API dei socket.}
\end{figure}

I socket sono l'elemento centrale per creare applicazioni di rete: realizzano, come una scatola nera,
le funzionalità di \textbf{trasporto} (TCP/UDP) e di \textbf{rete} (IP). Esistono librerie e
\emph{middleware} che ``avvolgono'' i socket semplificandone l'uso e offrendo funzioni più complesse, ma le
API di base sono ancora molto usate.

Le API per socket TCP e UDP esistono in quasi tutti i linguaggi e sistemi operativi (C, C++, C\#, Java,
Python, \dots). Nel mondo Unix si usano i \textbf{socket di Berkeley} (detti anche socket BSD o POSIX): API
native in C, le stesse che useremo noi.

\info{Tutto è un file}{
    I socket seguono la filosofia Unix per cui ``\emph{tutto è un file}'': un socket è identificato da un
    \textbf{file descriptor}, come un file aperto, e si può chiudere con \texttt{close()} o leggere con
    \texttt{read()}.
}

% ══════════════════════════════════════════════════════════════
\section{UDP e TCP: le due sequenze}
% ══════════════════════════════════════════════════════════════

\begin{figure}[H]
\centering
\begin{tikzpicture}[every node/.style={font=\scriptsize\sffamily}, st/.style={draw, rounded corners, fill=#1, minimum width=2.2cm, minimum height=0.6cm, align=center}]
  \begin{scope}
    \node[font=\small\bfseries\sffamily] at (1.4,1) {Senza connessione (UDP)};
    \node[font=\bfseries\sffamily] at (0,0.4) {Client}; \node[font=\bfseries\sffamily] at (2.8,0.4) {Server};
    \node[st=lapp!40] (c1) at (0,-0.2) {crea socket}; \node[st=lapp!40] (s1) at (2.8,-0.2) {crea socket};
    \node[st=ltra!40] (c3) at (0,-1.8) {invia messaggio}; \node[st=ltra!40] (s3) at (2.8,-1.8) {legge messaggio};
    \node[st=ltra!40] (c4) at (0,-2.6) {legge risposta}; \node[st=ltra!40] (s4) at (2.8,-2.6) {invia risposta};
    \node[st=lnet!40] (c5) at (0,-3.4) {chiude socket};
    \foreach \a/\b in {c1/c3, c3/c4, c4/c5, s1/s3, s3/s4} \draw[-{Stealth}] (\a) -- (\b);
    \draw[msg, netblue] (c3) -- (s3); \draw[msg, netgreen] (s4) -- (c4);
  \end{scope}
  \begin{scope}[xshift=9cm]
    \node[font=\small\bfseries\sffamily] at (1.4,1) {Orientata alla connessione (TCP)};
    \node[font=\bfseries\sffamily] at (0,0.4) {Client}; \node[font=\bfseries\sffamily] at (2.8,0.4) {Server};
    \node[st=lapp!40] (c1) at (0,-0.2) {crea socket}; \node[st=lapp!40] (s1) at (2.8,-0.2) {crea socket};
    \node[st=netred!15] (c2) at (0,-1.0) {connetti}; \node[st=netred!15] (s2) at (2.8,-1.0) {attende\\connessioni};
    \node[st=ltra!40] (c3) at (0,-1.8) {invia messaggio}; \node[st=ltra!40] (s3) at (2.8,-1.8) {legge messaggio};
    \node[st=ltra!40] (c4) at (0,-2.6) {legge risposta}; \node[st=ltra!40] (s4) at (2.8,-2.6) {invia risposta};
    \node[st=lnet!40] (c5) at (0,-3.4) {chiude socket}; \node[st=lnet!40] (s5) at (2.8,-3.4) {chiude socket\\del client};
    \foreach \a/\b in {c1/c2, c2/c3, c3/c4, c4/c5, s1/s2, s2/s3, s3/s4, s4/s5} \draw[-{Stealth}] (\a) -- (\b);
    \draw[msg, netred, {Stealth}-{Stealth}] (c2) -- (s2);
    \draw[msg, netblue] (c3) -- (s3); \draw[msg, netgreen] (s4) -- (c4);
  \end{scope}
  \foreach \y/\t in {-0.2/creazione, -1.0/connessione, -2.2/messaggi, -3.4/distruzione} \node[font=\scriptsize\itshape\sffamily, text=black!60] at (6.2,\y) {\t};
\end{tikzpicture}
\caption{Le fasi di una comunicazione senza connessione e con connessione.}
\end{figure}

% ══════════════════════════════════════════════════════════════
\section{Le strutture dati}
% ══════════════════════════════════════════════════════════════

I socket del dominio Internet in C usano alcune strutture per descrivere indirizzi e porte degli host:

\begin{lstlisting}[language={[POSIX]C}]
#include <netinet/in.h>
#include <arpa/inet.h>     // per inet_addr()
#include <sys/types.h>     // alcuni tipi personalizzati

struct sockaddr_in {
    short          sin_family;  // famiglia dell'indirizzo, di solito AF_INET (IPv4)
    unsigned short sin_port;    // numero di porta, per esempio htons(3490)
    struct in_addr sin_addr;    // indirizzo IP (vedi sotto)
    char           sin_zero[8]; // zeri: la struttura ha la stessa dimensione di sockaddr
};

struct in_addr {
    unsigned long s_addr;       // indirizzo IP: inet_addr("...") oppure INADDR_ANY
};
\end{lstlisting}

\begin{itemize}
    \item \texttt{sin\_zero} serve solo a rendere la struttura grande quanto la generica \texttt{struct sockaddr}:
          per questo alle funzioni si passa \texttt{(struct sockaddr *)\&indirizzo}.
    \item \texttt{inet\_addr("192.168.1.10")} converte un indirizzo in notazione puntata nel formato binario;
          \texttt{INADDR\_ANY} significa ``qualsiasi indirizzo di questa macchina''.
\end{itemize}

\warning{L'ordine dei byte: htons()}{
    Macchine diverse memorizzano i numeri su più byte in ordine diverso (\emph{little endian} o \emph{big endian}).
    In rete si usa per convenzione il \textbf{network byte order} (big endian). Per questo porte e indirizzi vanno
    convertiti: \texttt{htons()} (\emph{host to network short}) per le porte a 16 bit, \texttt{htonl()} per i
    valori a 32 bit, e \texttt{ntohs()}/\texttt{ntohl()} per il percorso inverso.
    \begin{center}
    \begin{tikzpicture}[every node/.style={font=\scriptsize\ttfamily}]
      \node[anchor=east, font=\scriptsize\sffamily] at (-0.2,0) {porta 8080 = 0x1F90, su x86:};
      \node[field, fill=lapp!50, minimum width=1cm] at (0.5,0) {90}; \node[field, fill=lnet!50, minimum width=1cm] at (1.5,0) {1F};
      \node[anchor=east, font=\scriptsize\sffamily] at (-0.2,-0.8) {dopo htons(8080), in rete:};
      \node[field, fill=lnet!50, minimum width=1cm] at (0.5,-0.8) {1F}; \node[field, fill=lapp!50, minimum width=1cm] at (1.5,-0.8) {90};
    \end{tikzpicture}
    \end{center}
}

% ══════════════════════════════════════════════════════════════
\section{Le funzioni di base}
% ══════════════════════════════════════════════════════════════

\subsection{Creazione del socket: socket()}

\begin{lstlisting}[language={[POSIX]C}]
#include <sys/socket.h>
int sockfd = socket(int domain, int type, int protocol);
\end{lstlisting}

\begin{itemize}
    \item \texttt{domain}: la famiglia di indirizzi, come nella struttura (\texttt{AF\_INET});
    \item \texttt{type}: il tipo di socket: \texttt{SOCK\_STREAM} $\rightarrow$ \textbf{TCP}, \texttt{SOCK\_DGRAM}
          $\rightarrow$ \textbf{UDP};
    \item \texttt{protocol}: un protocollo specifico da usare nel socket; di solito \texttt{0} (quello di default);
    \item \texttt{sockfd}: il \textbf{file descriptor} del nuovo socket (negativo in caso di errore).
\end{itemize}

\subsection{Associazione a indirizzo e porta: bind()}

\begin{lstlisting}[language={[POSIX]C}]
#include <sys/socket.h>
int val = bind(int socket, const struct sockaddr *address, socklen_t address_len);
\end{lstlisting}

\begin{itemize}
    \item \texttt{socket}: il file descriptor del socket da associare;
    \item \texttt{address}: puntatore a una \texttt{sockaddr} con porta e indirizzo su cui mettersi in ascolto;
          spesso \texttt{sin\_addr} è \texttt{INADDR\_ANY}, così si ascolta su tutti gli indirizzi della macchina;
    \item \texttt{address\_len}: la dimensione della struttura (con \texttt{sizeof()});
    \item \texttt{val}: 0 in caso di successo, $-1$ altrimenti.
\end{itemize}

\warning{bind() serve al server}{
    Il server \textbf{deve} associare il socket a una porta precisa (i client devono sapere dove trovarlo).
    Sul client \texttt{bind()} è \textbf{facoltativa}: se non la si chiama, il sistema operativo sceglie
    automaticamente una porta libera.
}

\subsection{Invio e ricezione in UDP: sendto() e recvfrom()}

\begin{lstlisting}[language={[POSIX]C}]
#include <sys/socket.h>
int ob = sendto(int osock, const void *obuf, size_t olen, int oflags,
                const struct sockaddr *oaddr, socklen_t oaddr_len);
int ib = recvfrom(int isock, void *ibuf, size_t ilen, int iflags,
                  struct sockaddr *iaddr, socklen_t *iaddr_len);
\end{lstlisting}

\begin{itemize}
    \item \texttt{osock}/\texttt{isock}: il socket su cui inviare/ricevere;
    \item \texttt{obuf}/\texttt{ibuf}: il buffer con il messaggio da inviare / dove mettere quello ricevuto;
    \item \texttt{olen}/\texttt{ilen}: la lunghezza in byte;
    \item \texttt{oflags}/\texttt{iflags}: opzioni, di solito \texttt{0} (oppure \texttt{MSG\_WAITALL} per aspettare
          tutti i byte richiesti);
    \item \texttt{oaddr}: l'indirizzo del \textbf{destinatario}; \texttt{iaddr}: dove la funzione \textbf{scrive}
          l'indirizzo del mittente;
    \item \texttt{ob}/\texttt{ib}: il numero di byte effettivamente inviati/ricevuti.
\end{itemize}

\info{Perché in UDP serve l'indirizzo ogni volta}{
    In UDP non c'è connessione: il socket non ``sa'' con chi sta parlando. Per questo \texttt{sendto()} vuole
    l'indirizzo del destinatario a ogni invio, e \texttt{recvfrom()} restituisce l'indirizzo di chi ha mandato il
    messaggio, così che il server possa rispondergli.
}

\subsection{Chiusura: close()}

\begin{lstlisting}[language={[POSIX]C}]
#include <unistd.h>
int val = close(int socket);   // 0 in caso di successo, -1 altrimenti
\end{lstlisting}

% ══════════════════════════════════════════════════════════════
\section{Esempio completo: client e server UDP}
% ══════════════════════════════════════════════════════════════

\begin{figure}[H]
\centering
\begin{tikzpicture}[yscale=0.7, every node/.style={font=\scriptsize\ttfamily}]
  \node[font=\footnotesize\bfseries\sffamily] at (0,0.6) {Client}; \node[font=\footnotesize\bfseries\sffamily] at (6,0.6) {Server};
  \draw[thick, black!50] (0,0) -- (0,-7.3); \draw[thick, black!50] (6,0) -- (6,-7.3);
  \node[fill=white, draw, rounded corners] at (0,-0.4) {socket()}; \node[fill=white, draw, rounded corners] at (6,-0.4) {socket()};
  \node[fill=white, draw, rounded corners] at (6,-1.4) {bind()};
  \node[fill=white, draw, rounded corners] (rf) at (6,-2.8) {recvfrom()};
  \node[fill=white, draw, rounded corners] (st) at (0,-2.4) {sendto()};
  \draw[msg, netblue] (st) -- node[above, sloped, font=\scriptsize\sffamily] {``Hello from client''} (rf);
  \node[fill=white, draw, rounded corners] (st2) at (6,-4.2) {sendto()};
  \node[fill=white, draw, rounded corners] (rf2) at (0,-4.8) {recvfrom()};
  \draw[msg, netgreen] (st2) -- node[above, sloped, font=\scriptsize\sffamily] {``Hello from server''} (rf2);
  \node[fill=white, draw, rounded corners] at (0,-6.5) {close()}; \node[fill=white, draw, rounded corners] at (6,-6.5) {close()};
  \node[anchor=west, font=\scriptsize\sffamily, text width=6cm, align=left] at (7,-2.8) {\textbf{recvfrom} si blocca finché non arriva un messaggio e restituisce IP e porta del client};
  \draw[{Stealth}-{Stealth}, dashed, black!50] (-0.8,-2.2) -- node[left, font=\scriptsize\sffamily, align=center] {si può\\ripetere} (-0.8,-5.0);
\end{tikzpicture}
\caption{La sequenza di chiamate per un semplice scambio UDP.}
\end{figure}

\begin{lstlisting}[language={[POSIX]C}, style=wnumbers, title={udp\_server.c}]
#include <stdio.h>
#include <stdlib.h>
#include <string.h>
#include <unistd.h>
#include <arpa/inet.h>
#include <sys/socket.h>

#define PORT    8080
#define MAXLINE 1024

int main(void) {
    int sockfd;
    char buffer[MAXLINE];
    const char *hello = "Hello from server";
    struct sockaddr_in servaddr, cliaddr;

    // 1. Creazione del socket UDP
    if ((sockfd = socket(AF_INET, SOCK_DGRAM, 0)) < 0) {
        perror("socket creation failed");
        exit(EXIT_FAILURE);
    }

    memset(&servaddr, 0, sizeof(servaddr));
    memset(&cliaddr, 0, sizeof(cliaddr));
    servaddr.sin_family      = AF_INET;
    servaddr.sin_addr.s_addr = INADDR_ANY;     // tutte le interfacce
    servaddr.sin_port        = htons(PORT);

    // 2. Associazione alla porta 8080
    if (bind(sockfd, (const struct sockaddr *)&servaddr, sizeof(servaddr)) < 0) {
        perror("bind failed");
        exit(EXIT_FAILURE);
    }

    // 3. Ricezione: cliaddr viene riempita con indirizzo e porta del client
    socklen_t len = sizeof(cliaddr);
    int n = recvfrom(sockfd, buffer, MAXLINE - 1, 0,
                     (struct sockaddr *)&cliaddr, &len);
    buffer[n] = '\0';
    printf("Ricevuto \"%s\"\n", buffer);

    // 4. Risposta proprio a chi ci ha scritto
    sendto(sockfd, hello, strlen(hello), 0, (const struct sockaddr *)&cliaddr, len);
    printf("Hello message sent.\n");

    close(sockfd);
    return 0;
}
\end{lstlisting}

\begin{lstlisting}[language={[POSIX]C}, style=wnumbers, title={udp\_client.c}]
#include <stdio.h>
#include <stdlib.h>
#include <string.h>
#include <unistd.h>
#include <arpa/inet.h>
#include <sys/socket.h>

#define PORT    8080
#define MAXLINE 1024

int main(void) {
    int sockfd;
    char buffer[MAXLINE];
    const char *hello = "Hello from client";
    struct sockaddr_in servaddr;

    if ((sockfd = socket(AF_INET, SOCK_DGRAM, 0)) < 0) {
        perror("socket creation failed");
        exit(EXIT_FAILURE);
    }

    memset(&servaddr, 0, sizeof(servaddr));
    servaddr.sin_family      = AF_INET;
    servaddr.sin_port        = htons(PORT);
    servaddr.sin_addr.s_addr = inet_addr("127.0.0.1");  // server sulla stessa macchina

    // Nessuna bind(): la porta del client la sceglie il sistema operativo
    sendto(sockfd, hello, strlen(hello), 0, (const struct sockaddr *)&servaddr, sizeof(servaddr));
    printf("Hello message sent.\n");

    socklen_t len = sizeof(servaddr);
    int n = recvfrom(sockfd, buffer, MAXLINE - 1, 0, (struct sockaddr *)&servaddr, &len);
    buffer[n] = '\0';
    printf("Ricevuto \"%s\"\n", buffer);

    close(sockfd);
    return 0;
}
\end{lstlisting}

\begin{lstlisting}[language=bash, title={Compilazione ed esecuzione, in due terminali}]
$ gcc -o udp_server udp_server.c
$ gcc -o udp_client udp_client.c
$ ./udp_server        # terminale 1: resta in attesa
$ ./udp_client        # terminale 2
\end{lstlisting}

\warning{Parlare con un altro computer}{
    Nelle trasparenze il client usa \texttt{INADDR\_ANY} anche come indirizzo di destinazione, il che funziona
    solo in locale. Per parlare con un'altra macchina si usa \texttt{inet\_addr()} con il suo indirizzo, per
    esempio \texttt{inet\_addr("192.168.1.10")}. Attenzione anche a lasciare un byte libero nel buffer per il
    terminatore \texttt{'\textbackslash 0'}: se si riceve esattamente \texttt{MAXLINE} byte,
    \texttt{buffer[n]} andrebbe fuori dall'array.
}

\subsection{Riassunto della sequenza UDP}

\begin{itemize}
    \item Client e server aprono entrambi un socket.
    \item Il server lo associa a una porta precisa con \texttt{bind()}.
    \item Si scambiano messaggi con \texttt{sendto()} e \texttt{recvfrom()}; IP e porta del client si ricavano
          da \texttt{recvfrom()}. Si può ciclare su \texttt{sendto}/\texttt{recvfrom} per proseguire la conversazione.
    \item Alla fine entrambi chiudono il socket. Se chiude solo il client, il server può continuare ad accettare
          messaggi da altri client.
\end{itemize}

% ══════════════════════════════════════════════════════════════
\section{Da UDP a TCP}
% ══════════════════════════════════════════════════════════════

A differenza di UDP, in TCP client e server devono fare l'\textbf{handshake} prima di scambiarsi messaggi.
Sul server servono \textbf{due} socket:

\begin{itemize}
    \item un \textbf{socket di benvenuto} (\emph{welcoming socket}), sempre attivo, su cui i client fanno
          l'handshake con il server;
    \item un \textbf{socket dedicato al client}, creato dopo l'handshake, che client e server usano per comunicare.
\end{itemize}

\begin{figure}[H]
\centering
\begin{tikzpicture}[every node/.style={font=\footnotesize\sffamily}]
  \node (c1) at (0,1.5) {\icon[0.8cm]{pc}}; \node[left=0pt of c1] {client 1};
  \node (c2) at (0,0) {\icon[0.9cm]{laptop}}; \node[left=0pt of c2] {client 2};
  \node (c3) at (0,-1.5) {\icon[0.8cm]{pc}}; \node[left=0pt of c3] {client 3};
  \node[draw, very thick, netred, fill=netred!10, rounded corners, minimum width=1.6cm, minimum height=0.8cm, align=center] (w) at (4.5,0) {socket di\\benvenuto};
  \node[draw, fill=llink!60, minimum height=0.6cm, minimum width=1.8cm] (q) at (7,0) {coda: 3 2 1};
  \node[draw=netblue, thick, rounded corners, minimum width=6.8cm, minimum height=4.3cm] (srv) at (7.6,0) {};
  \node[anchor=north east] at (srv.north east) {\icon[0.5cm]{server}\ server};
  \node[draw, fill=ltra!50, rounded corners, align=center] (a) at (9.6,0.9) {socket dedicato\\al client 1};
  \node[anchor=west, font=\ttfamily\scriptsize] at (4.0,1.4) {listen()};
  \node[anchor=west, font=\ttfamily\scriptsize] at (8.1,-0.9) {accept()};
  \foreach \c in {c1,c2,c3} \draw[msg, netred] (\c) -- node[above, sloped, font=\ttfamily\scriptsize] {connect()} (w);
  \draw[-{Stealth}, thick] (w) -- (q); \draw[-{Stealth}, thick] (q) -- (a);
  \draw[msg, netgreen, {Stealth}-{Stealth}, very thick] (c1) to[bend left=25] node[above, font=\scriptsize\sffamily] {comunicazione} (a);
\end{tikzpicture}
\caption{\texttt{listen()} apre una coda in cui si accumulano le richieste di connessione; \texttt{accept()} ne estrae una e crea il socket dedicato a quel client.}
\end{figure}

Il three-way handshake avviene dentro il livello di trasporto ed è \textbf{del tutto invisibile} ai
programmi client e server: una volta accettata la connessione, la comunicazione si sposta sul socket
appena creato.

\subsection{connect() lato client}

\begin{lstlisting}[language={[POSIX]C}]
#include <sys/socket.h>
int val = connect(int socket, const struct sockaddr *address, socklen_t address_len);
\end{lstlisting}

\begin{itemize}
    \item \texttt{socket}: il socket con cui connettersi;
    \item \texttt{address}: l'indirizzo del \textbf{server};
    \item \texttt{val}: 0 in caso di successo, $-1$ altrimenti (per esempio se nessuno è in ascolto).
\end{itemize}

\subsection{listen() e accept() lato server}

\begin{lstlisting}[language={[POSIX]C}]
#include <sys/socket.h>
int val        = listen(int socket, int backlog);
int new_sockfd = accept(int socket, struct sockaddr *address, socklen_t *address_len);
\end{lstlisting}

\begin{itemize}
    \item \texttt{socket}: il socket su cui attendere le connessioni (deve essere già associato con \texttt{bind()});
    \item \texttt{backlog}: la lunghezza massima della \textbf{coda} delle connessioni in attesa;
    \item \texttt{address}: un puntatore nullo, oppure una \texttt{sockaddr} in cui \texttt{accept()} scrive
          l'indirizzo del client che si è connesso;
    \item \texttt{new\_sockfd}: il \textbf{nuovo socket} su cui client e server comunicheranno.
\end{itemize}

\texttt{accept()} è \textbf{bloccante}: il server resta fermo lì finché un client non si connette.

\subsection{send() e read() in TCP}

\begin{lstlisting}[language={[POSIX]C}]
#include <sys/socket.h>
int ob = send(int osock, const void *obuf, size_t olen, int flags);
#include <unistd.h>
int ib = read(int isock, void *ibuf, size_t ilen);
\end{lstlisting}

Sono più semplici delle versioni UDP: gli indirizzi di client e server sono già stati fissati durante la
connessione, e non serve ripeterli. \texttt{flags} vale di solito 0; \texttt{ob} e \texttt{ib} sono i byte
effettivamente inviati e ricevuti (\texttt{read()} restituisce 0 quando l'altro lato ha chiuso la connessione).

\warning{TCP è un flusso di byte, non di messaggi}{
    Con TCP non esistono ``messaggi'': una \texttt{send()} di 100 byte può arrivare all'altra parte in due
    \texttt{read()} da 60 e 40 byte, e due \texttt{send()} possono arrivare in una sola \texttt{read()}. Un
    protocollo applicativo su TCP deve quindi stabilire come delimitare i messaggi (una lunghezza in testa, un
    carattere terminatore come \texttt{\textbackslash r\textbackslash n} in HTTP, \dots).
}

% ══════════════════════════════════════════════════════════════
\section{Esempio completo: client e server TCP}
% ══════════════════════════════════════════════════════════════

\begin{lstlisting}[language={[POSIX]C}, style=wnumbers, title={tcp\_server.c}]
#include <stdio.h>
#include <stdlib.h>
#include <string.h>
#include <unistd.h>
#include <arpa/inet.h>
#include <sys/socket.h>

#define PORT 8080

int main(void) {
    int sockfd, new_socket;
    char buffer[1024];
    const char *hello = "Hello from server";
    struct sockaddr_in servaddr, cliaddr;

    // 1. Socket di benvenuto (TCP)
    if ((sockfd = socket(AF_INET, SOCK_STREAM, 0)) < 0) {
        perror("socket creation failed");
        exit(EXIT_FAILURE);
    }
    memset(&servaddr, 0, sizeof(servaddr));
    memset(&cliaddr, 0, sizeof(cliaddr));
    servaddr.sin_family      = AF_INET;
    servaddr.sin_addr.s_addr = INADDR_ANY;
    servaddr.sin_port        = htons(PORT);

    // 2. bind + listen: coda di al massimo 3 connessioni in attesa
    if (bind(sockfd, (const struct sockaddr *)&servaddr, sizeof(servaddr)) < 0) {
        perror("bind failed");
        exit(EXIT_FAILURE);
    }
    if (listen(sockfd, 3) < 0) {
        perror("listen");
        exit(EXIT_FAILURE);
    }

    // 3. accept: si blocca fino alla connessione di un client
    socklen_t addrlen = sizeof(cliaddr);
    if ((new_socket = accept(sockfd, (struct sockaddr *)&cliaddr, &addrlen)) < 0) {
        perror("accept");
        exit(EXIT_FAILURE);
    }

    // 4. Dialogo sul socket dedicato
    int n = read(new_socket, buffer, sizeof(buffer) - 1);
    buffer[n > 0 ? n : 0] = '\0';
    printf("Ricevuto \"%s\"\n", buffer);
    send(new_socket, hello, strlen(hello), 0);
    printf("Hello message sent.\n");

    // 5. Chiusura: prima il socket del client, poi quello di benvenuto
    close(new_socket);
    close(sockfd);
    return 0;
}
\end{lstlisting}

\begin{lstlisting}[language={[POSIX]C}, style=wnumbers, title={tcp\_client.c}]
#include <stdio.h>
#include <stdlib.h>
#include <string.h>
#include <unistd.h>
#include <arpa/inet.h>
#include <sys/socket.h>

#define PORT 8080

int main(void) {
    int sockfd;
    char buffer[1024];
    const char *hello = "Hello from client";
    struct sockaddr_in servaddr;

    if ((sockfd = socket(AF_INET, SOCK_STREAM, 0)) < 0) {
        perror("socket creation failed");
        exit(EXIT_FAILURE);
    }
    memset(&servaddr, 0, sizeof(servaddr));
    servaddr.sin_family      = AF_INET;
    servaddr.sin_port        = htons(PORT);
    servaddr.sin_addr.s_addr = inet_addr("127.0.0.1");

    // L'handshake a tre vie avviene qui dentro, invisibile al programma
    if (connect(sockfd, (struct sockaddr *)&servaddr, sizeof(servaddr)) < 0) {
        printf("\nConnection Failed\n");
        return -1;
    }

    send(sockfd, hello, strlen(hello), 0);
    printf("Hello message sent.\n");

    int n = read(sockfd, buffer, sizeof(buffer) - 1);
    buffer[n > 0 ? n : 0] = '\0';
    printf("Ricevuto \"%s\"\n", buffer);

    close(sockfd);
    return 0;
}
\end{lstlisting}

\subsection{Riassunto della sequenza TCP}

\begin{figure}[H]
\centering
\begin{tikzpicture}[yscale=0.66, every node/.style={font=\scriptsize\ttfamily}]
  \node[font=\footnotesize\bfseries\sffamily] at (0,0.7) {Client}; \node[font=\footnotesize\bfseries\sffamily] at (6,0.7) {Server};
  \draw[thick, black!50] (0,0) -- (0,-10.2); \draw[thick, black!50] (6,0) -- (6,-10.2);
  \foreach \x/\y/\t in {0/-0.3/socket(), 6/-0.3/socket(), 6/-1.1/bind(), 6/-1.9/listen(), 6/-2.9/accept(), 0/-3.9/connect(), 0/-5.3/send(), 6/-5.9/read(), 6/-6.7/send(), 0/-7.3/read(), 0/-8.3/close(), 6/-8.9/read(), 6/-9.7/close()}
    \node[fill=white, draw, rounded corners] (\t\y) at (\x,\y) {\t};
  \draw[msg, netred, {Stealth}-{Stealth}] (0.7,-3.9) -- node[above, font=\scriptsize\sffamily] {handshake} (5.3,-3.4);
  \node[anchor=west, font=\scriptsize\sffamily, text=netred] at (6.9,-2.9) {bloccata fino al nuovo client};
  \draw[msg, netblue] (0.6,-5.3) -- (5.4,-5.9);
  \draw[msg, netgreen] (5.4,-6.7) -- (0.6,-7.3);
  \draw[msg, black!60, {Stealth}-{Stealth}] (0.7,-8.3) -- node[above, font=\scriptsize\sffamily] {rilascio} (5.3,-8.9);
  \node[anchor=west, font=\scriptsize\sffamily, text width=5cm, align=left] at (6.9,-8.9) {la \texttt{read()} restituisce 0: il client ha chiuso};
  \draw[-{Stealth}, dashed, netorange, thick] (7.0,-9.7) to[out=0, in=0] node[right, font=\scriptsize\sffamily, text width=3cm, align=left] {può tornare ad \texttt{accept()} per un nuovo client} (6.9,-2.9);
\end{tikzpicture}
\caption{La sequenza di chiamate TCP.}
\end{figure}

\begin{itemize}
    \item Client e server aprono entrambi un socket.
    \item Il server lo associa a una porta (\texttt{bind}), si mette in ascolto (\texttt{listen}) e accetta le connessioni
          (\texttt{accept}); il client si connette (\texttt{connect}).
    \item Si scambiano messaggi con \texttt{send} e \texttt{read}; IP e porta del client si ricavano dal socket
          appena creato. Si può ciclare su \texttt{send}/\texttt{read}.
    \item Alla fine entrambi chiudono il socket. Il server può tornare ad \texttt{accept()} per accogliere un altro client.
\end{itemize}

\warning{La read() prima della close() e i quattro minuti di attesa}{
    Il server dovrebbe fare un'ulteriore \texttt{read()} prima di chiudere il socket del client, così da
    completare il \textbf{rilascio} della connessione avviato dal client; altrimenti la porta resta occupata
    (stato \texttt{TIME\_WAIT}) e bisogna aspettare circa \textbf{4 minuti} prima di poterla riusare. Durante le
    prove si può evitare il problema impostando sul socket di benvenuto l'opzione \texttt{SO\_REUSEADDR}:
    \texttt{int on = 1; setsockopt(sockfd, SOL\_SOCKET, SO\_REUSEADDR, \&on, sizeof(on));} prima di \texttt{bind()}.
}

% ══════════════════════════════════════════════════════════════
\section{La traduzione DNS in C}
% ══════════════════════════════════════════════════════════════

Per creare un socket verso un server serve il suo \textbf{indirizzo IP}. Ma su Internet gli host si
identificano con i \textbf{nomi}, e i socket di Berkeley lavorano \textbf{solo con indirizzi IP}. Non è
strano: i socket lavorano al livello di trasporto, i nomi vivono al livello applicativo. Per connetterci a un
host dato il nome dobbiamo quindi usare il DNS e ricavarne l'indirizzo.

\subsection{La struttura hostent}

La funzione \texttt{gethostbyname()} contatta i server DNS e restituisce una struttura \texttt{hostent} con
alcune informazioni sull'host: nome canonico, eventuali alias, uno o più indirizzi.

\begin{lstlisting}[language={[POSIX]C}]
#include <netdb.h>
struct hostent {
    char  *h_name;       // nome originale (canonico) dell'host
    char **h_aliases;    // lista di alias (terminata da NULL)
    int    h_addrtype;   // famiglia dell'indirizzo (di solito AF_INET)
    int    h_length;     // lunghezza dell'indirizzo (di solito 4 byte)
    char **h_addr_list;  // lista di indirizzi (terminata da NULL)
};
#define h_addr h_addr_list[0]   // per compatibilita': il primo indirizzo
\end{lstlisting}

Il DNS può restituire più indirizzi (in ordine casuale, se si interroga un server round-robin); di solito si
usa il primo, \texttt{h\_addr\_list[0]} (o \texttt{h\_addr}). Gli indirizzi restituiti si possono convertire in
\texttt{struct in\_addr}, la stessa usata dai socket.

\subsection{gethostbyname()}

\begin{lstlisting}[language={[POSIX]C}]
#include <netdb.h>
struct hostent *host_info = gethostbyname(const char *name);   // NULL in caso di errore
\end{lstlisting}

Esiste anche \texttt{gethostbyaddr()}, che restituisce una \texttt{hostent} a partire da un indirizzo IP (ricerca inversa).

\begin{lstlisting}[language={[POSIX]C}, style=wnumbers, title={dns\_lookup.c}]
#include <stdio.h>
#include <netdb.h>
#include <arpa/inet.h>

int main(int argc, char **argv) {
    if (argc < 2) {
        printf("No hostname specified\n");
        return 1;
    }

    // Interroga il DNS
    struct hostent *info = gethostbyname(argv[1]);
    if (info == NULL) {
        printf("Could not resolve hostname %s :(\n", argv[1]);
        return 1;
    }

    printf("Canonical name:\n\t%s\n", info->h_name);

    printf("Aliases:\n");
    for (int i = 0; info->h_aliases[i] != NULL; i++)
        printf("\t%s\n", info->h_aliases[i]);

    printf("IPs:\n");
    for (int i = 0; info->h_addr_list[i] != NULL; i++)
        printf("\t%s\n", inet_ntoa(*(struct in_addr *)info->h_addr_list[i]));

    return 0;
}
\end{lstlisting}

\warning{Occhio agli errori nelle trasparenze}{
    Il codice delle trasparenze (in C++) contiene alcune sviste da non copiare: controlla
    \texttt{server\_name == NULL} invece del risultato di \texttt{gethostbyname()}, e nel ciclo degli alias
    l'incremento \texttt{i++} è fuori dalle graffe, quindi il ciclo non terminerebbe mai. La versione qui sopra le
    corregge. Notate anche che \texttt{gethostbyname()} è considerata obsoleta: i programmi moderni usano
    \texttt{getaddrinfo()}, che supporta anche IPv6.
}

% ══════════════════════════════════════════════════════════════
\section{Una richiesta HTTP scritta a mano}
% ══════════════════════════════════════════════════════════════

Mettiamo insieme tutto: un programma che, da zero, invia una richiesta HTTP \texttt{HEAD} per la pagina
``cenni storici'' del sito dell'Ateneo (\url{http://www.unina.it/chi-siamo/cenni-storici}). Servono tre
elementi: il \textbf{nome del server}, la \textbf{porta} (80) e il \textbf{percorso} della risorsa.

\begin{lstlisting}[language={[POSIX]C}, style=wnumbers, title={http\_head.c}]
#include <stdio.h>
#include <string.h>
#include <unistd.h>
#include <netdb.h>
#include <arpa/inet.h>
#include <sys/socket.h>

int main(void) {
    // 1. Socket TCP
    int sd = socket(AF_INET, SOCK_STREAM, 0);
    if (sd < 0) { perror("socket"); return 1; }

    // 2. DNS: dal nome all'indirizzo IP
    struct hostent *server = gethostbyname("www.unina.it");
    if (server == NULL) {
        printf("Could not resolve hostname :(\n");
        close(sd);
        return 1;
    }

    struct sockaddr_in serv_addr;
    memset(&serv_addr, 0, sizeof(serv_addr));
    serv_addr.sin_family = AF_INET;
    serv_addr.sin_port   = htons(80);                       // porta di HTTP
    memcpy(&serv_addr.sin_addr.s_addr, server->h_addr, server->h_length);

    // 3. Connessione TCP al server Web
    if (connect(sd, (struct sockaddr *)&serv_addr, sizeof(serv_addr)) < 0) {
        perror("connect");
        close(sd);
        return 1;
    }

    // 4. La richiesta HTTP: testo, righe chiuse da \r\n, riga vuota finale
    const char *request = "HEAD /chi-siamo/cenni-storici HTTP/1.1\r\n"
                          "Host: www.unina.it\r\n"
                          "Connection: close\r\n"
                          "\r\n";
    if (send(sd, request, strlen(request), 0) < 0) {
        perror("send");
        close(sd);
        return 1;
    }
    printf("message sent:\n%s", request);

    // 5. La risposta
    char buffer[4096];
    int n = recv(sd, buffer, sizeof(buffer) - 1, 0);
    if (n > 0) {
        buffer[n] = '\0';
        printf("received %d bytes:\n%s", n, buffer);
    }

    close(sd);
    return 0;
}
\end{lstlisting}

\begin{itemize}
    \item \texttt{gethostbyname()} invoca il DNS e restituisce l'indirizzo del server nella \texttt{hostent}; lo si
          copia nella \texttt{serv\_addr} per la \texttt{connect()}.
    \item La richiesta è scritta nella stringa \texttt{request}; con \texttt{Connection: close} si chiede una
          connessione \textbf{non persistente}, perché ci interessa una sola pagina.
    \item Si notino i \texttt{\textbackslash r\textbackslash n} alla fine di ogni riga e la \textbf{riga vuota} finale che
          chiude gli header: senza, il server resterebbe in attesa del resto della richiesta.
\end{itemize}

\info{L'incapsulamento, toccato con mano}{
    Scrivendo una richiesta HTTP da zero si vede concretamente che un messaggio di livello applicativo viene
    inviato come \textbf{payload} di un messaggio di trasporto: per TCP la stringa \texttt{request} è solo una
    sequenza di byte, del tutto irrilevante nel contenuto.
}

% ══════════════════════════════════════════════════════════════
\section{Riepilogo}
% ══════════════════════════════════════════════════════════════

\riepilogo{
\begin{itemize}
    \item I socket (Berkeley/POSIX) sono l'API verso trasporto e rete, implementati dal sistema operativo; un
          socket è un file descriptor.
    \item Indirizzi in \texttt{struct sockaddr\_in}; porte e indirizzi in \textbf{network byte order} (\texttt{htons}).
    \item \textbf{UDP}: \texttt{socket}, \texttt{bind} (server), \texttt{sendto}/\texttt{recvfrom}, \texttt{close}.
    \item \textbf{TCP}: server \texttt{socket}, \texttt{bind}, \texttt{listen}, \texttt{accept} (nuovo socket dedicato);
          client \texttt{socket}, \texttt{connect}; poi \texttt{send}/\texttt{read} e \texttt{close}.
    \item Per passare da nome a indirizzo si usa il DNS: \texttt{gethostbyname()} e la struttura \texttt{hostent}.
\end{itemize}
}
